\chapter{Mạnh và Yếu}
\markboth{Mạnh và Yếu}{Strong and Weak}

\section{Một người tỏ ra rất mạnh nhưng không dám nói mình đang mệt.}
\begin{truyen}{Một người tỏ ra rất mạnh nhưng không dám nói mình đang mệt.}{Pretended strength can hide a quiet break.}
\chuhoa{A}{nh quen làm chỗ dựa nên không cho phép mình than. Càng lâu, mọi người càng tưởng anh không biết mệt, còn anh thì cạn dần ở bên trong. Có kiểu mạnh khiến người ta được nể, nhưng cũng có kiểu mạnh giả khiến người ta tự làm gãy mình.}
\textit{(He was used to being the support and no longer allowed himself to complain. Over time, people assumed he never tired, while he slowly emptied inside. Some strength earns respect, but some fake strength quietly breaks the person carrying it.)}
\end{truyen}
\begin{baihoc}
Mạnh không phải là không bao giờ yếu; mạnh là biết đối diện phần yếu một cách thành thật.
\textit{(Strength is not never being weak; it is facing your weakness honestly.)}
\end{baihoc}
\begin{ghinhoanh}
\textbf{Short English to remember:}

Real strength can speak of weakness.
\end{ghinhoanh}
\begin{tuhocnhanh}
1. Em có đang cố tỏ ra ổn trong khi thật ra rất mệt không? \textit{(Are you acting okay while actually very tired?)}

2. Em tin ai đủ để nói thật về trạng thái của mình? \textit{(Whom do you trust enough to tell the truth about your condition?)}
\end{tuhocnhanh}
\ngancach

\section{Một học sinh yếu nền nhưng không yếu ý chí.}
\begin{truyen}{Một học sinh yếu nền nhưng không yếu ý chí.}{Weak at first does not mean weak forever.}
\chuhoa{E}{m vào lớp với kiến thức hổng nhiều, làm bài liên tục thấp điểm. Nhưng em không bỏ, vẫn hỏi, vẫn luyện, vẫn chấp nhận học lại từ gốc. Cái yếu ban đầu không quyết định số phận, nếu bên trong còn một ý chí chịu lớn lên.}
\textit{(The student entered class with many learning gaps and consistently low scores. Yet there was no giving up, only questions, practice, and willingness to relearn from the roots. Initial weakness does not decide destiny if an inner will remains willing to grow.)}
\end{truyen}
\begin{baihoc}
Yếu ở điểm xuất phát chưa chắc yếu ở đích đến.
\textit{(Being weak at the starting point does not mean being weak at the finish.)}
\end{baihoc}
\begin{ghinhoanh}
\textbf{Short English to remember:}

Weakness can be temporary.
\end{ghinhoanh}
\begin{tuhocnhanh}
1. Điểm yếu nào của em thực ra vẫn có thể sửa được? \textit{(Which weakness of yours can still be repaired?)}

2. Em có đang dùng cái yếu hiện tại để tự kết luận quá sớm về mình không? \textit{(Are you using your current weakness to judge yourself too early?)}
\end{tuhocnhanh}
\ngancach

\section{Người mạnh nhất trong nhóm lại là người biết nhường.}
\begin{truyen}{Người mạnh nhất trong nhóm lại là người biết nhường.}{Gentleness can be a form of strength.}
\chuhoa{T}{rong tranh luận, ai cũng muốn phần mình thắng. Chỉ có một bạn biết dừng lại, nghe hết ý người khác, rồi nói ngắn gọn để kéo cuộc nói chuyện về đúng việc. Sức mạnh không phải lúc nào cũng ở giọng lớn; nhiều khi nó nằm ở khả năng tự chủ.}
\textit{(In the argument, everyone wanted to win. Only one student paused, listened fully, and then spoke briefly to bring the discussion back to the real issue. Strength is not always found in volume; often it lies in self-control.)}
\end{truyen}
\begin{baihoc}
Biết nhường không phải là yếu; nhiều khi đó là cách mạnh nhất để giữ cuộc việc khỏi đổ vỡ.
\textit{(Knowing how to yield is not weakness; often it is the strongest way to keep things from breaking apart.)}
\end{baihoc}
\begin{ghinhoanh}
\textbf{Short English to remember:}

Gentleness can be strong.
\end{ghinhoanh}
\begin{tuhocnhanh}
1. Em thường nghĩ mạnh là thắng người khác hay thắng chính mình? \textit{(Do you think strength is defeating others or mastering yourself?)}

2. Em có thể mềm lại ở đâu mà không hề mất giá trị? \textit{(Where can you become softer without losing value?)}
\end{tuhocnhanh}
\ngancach

\section{Một người yếu lòng nên dễ tin lời ngon ngọt.}
\begin{truyen}{Một người yếu lòng nên dễ tin lời ngon ngọt.}{Emotional weakness invites easy traps.}
\chuhoa{C}{hỉ cần ai đó nói vài câu nâng đỡ đúng lúc, cô đã thấy mình được hiểu và vội đặt niềm tin. Sau vài lần bị lợi dụng, cô mới hiểu yếu lòng không phải là có nhiều cảm xúc, mà là không đủ tỉnh để phân biệt ai thật ai không.}
\textit{(A few comforting words at the right time were enough for her to feel understood and quickly trust. After being used several times, she realized that being emotionally weak is not about having many feelings, but about lacking the clarity to distinguish sincerity from manipulation.)}
\end{truyen}
\begin{baihoc}
Một trái tim ấm vẫn cần một cái đầu tỉnh.
\textit{(A warm heart still needs a clear mind.)}
\end{baihoc}
\begin{ghinhoanh}
\textbf{Short English to remember:}

A warm heart still needs clarity.
\end{ghinhoanh}
\begin{tuhocnhanh}
1. Khi cô đơn, em có dễ tin nhầm người không? \textit{(When lonely, do you become too quick to trust the wrong person?)}

2. Em cần thêm sự tỉnh táo ở mối quan hệ nào? \textit{(In which relationship do you need more clarity?)}
\end{tuhocnhanh}
\ngancach

\section{Nhận ra lúc mình yếu để biết tìm sự giúp đỡ.}
\begin{truyen}{Nhận ra lúc mình yếu để biết tìm sự giúp đỡ.}{Knowing your weakness is a strength.}
\chuhoa{K}{hi cơn lo âu kéo dài, em không còn học nổi như trước. Thay vì tiếp tục giả vờ ổn, em tìm đến cô giáo chủ nhiệm và một người thân để nói thật. Chính giây phút dám nhận mình đang yếu lại mở ra con đường hồi phục.}
\textit{(When anxiety lasted too long, the student could no longer study as before. Instead of pretending to be fine, the truth was shared with a homeroom teacher and a family member. The moment of admitting weakness opened the road to healing.)}
\end{truyen}
\begin{baihoc}
Người mạnh không phải là người không cần ai, mà là người biết lúc nào cần ai.
\textit{(A strong person is not one who needs no one, but one who knows when someone is needed.)}
\end{baihoc}
\begin{ghinhoanh}
\textbf{Short English to remember:}

Admitting weakness can begin healing.
\end{ghinhoanh}
\begin{tuhocnhanh}
1. Điều gì em đang cố chịu một mình trong khi nên tìm người hỗ trợ? \textit{(What are you trying to bear alone even though you need support?)}

2. Ai là người em có thể tìm đến khi mình không ổn? \textit{(Who can you turn to when you are not okay?)}
\end{tuhocnhanh}
\ngancach
